%% kapa: exam=1 round=35 year=2024 period=1 form=A subject=real_estate subject_ko=부동산학원론 seq=30 number=110
%% source: Q-Net 감정평가사 기출문제 (cst00309) artlSeq=5213229 file=raw/1008/5213229/1교시 문제지 원본.pdf
%% conversion: kordoc 4.24.51 -> tools/gen1.py (auto); review=proofread-sonnet; flags=none
\begin{question}{110}
부동산금융에 관한 설명으로 옳은 것은? (단, 주어진 조건에 한함)
\begin{choices}
  \item 콜옵션(call option)은 저당대출 대출자에게 주어진 조기상환권이다.
  \item 금융기관은 위험을 줄이기 위해 부채감당률이 1보다 작은 대출안의 작은 순서대로 대출을 실행한다.
  \item 대출수수료와 조기상환수수료를 차입자가 부담하는 경우, 차입자의 실효이자율은 조기상환시점이 앞당겨 질수록 하락한다.
  \item 대출조건이 동일할 경우 대출채권의 듀레이션(평균회수기간)은 원리금균등분할상환방식이 원금균등분할상환방식보다 더 길다.
  \item 고정금리방식의 대출에서 총상환액은 원리금균등분할상환방식이 원금균등분할상환방식보다 더 작다.
\end{choices}
\end{question}
\answer{4}
